\documentclass[11pt,a4paper,fontset=windows]{ctexart}
% ===== 公共导言区：高等代数（一）期中模拟卷 =====
\usepackage{amsmath,amssymb}
\usepackage{unicode-math}
\setmathfont{TeX Gyre Termes Math}
\usepackage[margin=2.05cm,top=2.35cm,bottom=2.0cm,headsep=11pt,includefoot]{geometry}
\usepackage{fancyhdr}
\usepackage{lastpage}
\usepackage{enumitem}
\usepackage{array}
\usepackage{bm}
\usepackage{xcolor}
\usepackage{zhnumber}
\usepackage{needspace}

\definecolor{hdgray}{gray}{0.32}
\definecolor{lngray}{gray}{0.5}
\definecolor{notebg}{gray}{0.955}

\pagestyle{fancy}
\fancyhf{}
\renewcommand{\headrulewidth}{0.4pt}
\renewcommand{\footrulewidth}{0pt}
\fancyhead[L]{\small\color{hdgray}\songti 高等代数（一）\ 期中模拟试卷（\thepp）}
\fancyhead[R]{\small\color{hdgray}\songti 考试时间：120 分钟\quad 满分：100 分}
\fancyfoot[C]{\small\color{lngray}\songti 第 \thepage\ 页\quad 共 \pageref{LastPage}\ 页}

\setlength{\parindent}{0pt}
\setlength{\parskip}{2pt}
\linespread{1.16}

% 题目列表 1—8 连续编号
\newlist{qs}{enumerate}{1}
\setlist[qs]{label=\textbf{\arabic*.},ref=\arabic*,leftmargin=2.55em,
  labelsep=0.55em,itemsep=23pt,topsep=6pt,parsep=0pt}
\newlist{note}{enumerate}{1}
\setlist[note]{label=(\arabic*),leftmargin=1.75em,itemsep=1.5pt,topsep=2.5pt,parsep=0pt}
\newlist{sub}{enumerate}{1}
\setlist[sub]{label=(\arabic*),leftmargin=1.9em,itemsep=3pt,topsep=3pt,parsep=0pt}

\newcommand{\secA}[1]{\Needspace{9\baselineskip}\par\vspace{10pt}{\large\bfseries\songti #1}\par\vspace{3pt}\nopagebreak}
\newcommand{\blank}[1][2.6cm]{\underline{\hspace{#1}}}
\newcounter{pp}
\renewcommand{\thepp}{\chinese{pp}}

\setcounter{pp}{3}
\begin{document}

\begin{center}
{\LARGE\bfseries\songti 高等代数（一）期中模拟试卷（三）}\\[7pt]
{\normalsize\songti 适用范围：第一章\ 多项式（\S1\,--\,\S9）；第二章\ 行列式（\S2.1\,--\,\S2.7）}\\[3pt]
{\small\color{hdgray}\songti 多元多项式、对称多项式不在考查范围}
\end{center}
\vspace{-4pt}

\begin{center}
\begin{tabular}{@{}p{0.335\textwidth}p{0.335\textwidth}p{0.26\textwidth}@{}}
\songti 姓名：\blank[2.5cm] & \songti 学号：\blank[2.5cm] & \songti 得分：\blank[1.9cm]
\end{tabular}
\end{center}

\vspace{2pt}
\hrule height 0.5pt
\begin{center}
\setlength{\fboxsep}{6pt}%
\colorbox{notebg}{\parbox{0.94\textwidth}{\small\songti
\textbf{考生须知：}
\begin{note}
\item 本试卷共 8 题，满分 100 分，考试时间 120 分钟。
\item 计算题须写出必要的计算过程，证明题须给出完整的推理步骤；只写答案不给分。
\item 所有答案一律写在答题纸上，写在试卷上无效。
\end{note}}}
\end{center}
\hrule height 0.5pt
\vspace{4pt}

\secA{一、计算题（每小题 10 分，共 40 分）}
\begin{qs}
\item 设 $a,b$ 为常数。计算 $n$ 阶行列式
\[
D_n=\begin{vmatrix}
a+b&ab&0&\cdots&0\\
1&a+b&ab&\cdots&0\\
0&1&a+b&\ddots&\vdots\\
\vdots&\ddots&\ddots&\ddots&ab\\
0&\cdots&0&1&a+b
\end{vmatrix},
\]
并就 $a\neq b$ 与 $a=b$ 两种情形分别写出结果。

\item 计算 $n$ 阶行列式
\[
D_n=\begin{vmatrix}
1&2&3&\cdots&n\\
1&2^{2}&3^{2}&\cdots&n^{2}\\
1&2^{3}&3^{3}&\cdots&n^{3}\\
\vdots&\vdots&\vdots&&\vdots\\
1&2^{\,n}&3^{\,n}&\cdots&n^{\,n}
\end{vmatrix}.
\]

\item 用克拉默法则解下列线性方程组：
\[
\begin{cases}
2x_1+x_2-5x_3+x_4=-9,\\[2pt]
x_1-3x_2-6x_4=4,\\[2pt]
2x_2-x_3+2x_4=-4,\\[2pt]
x_1+4x_2-7x_3+6x_4=-17.
\end{cases}
\]

\item 设 $f(x)=x^{4}+1$。
\begin{sub}
\item 证明 $f(x)$ 在 $\mathbb{Q}[x]$ 中不可约；
\item 写出 $f(x)$ 在 $\mathbb{R}[x]$ 与 $\mathbb{C}[x]$ 中的标准分解式。
\end{sub}
\end{qs}

\secA{二、证明题（每小题 15 分，共 60 分）}
\begin{qs}[resume*]
\item 设 $f(x)\in P[x]$，$n$ 为正整数。证明：
\[
(x-1)\mid f(x)\iff \bigl(x^{n}-1\bigr)\mid f(x^{n}).
\]

\item （高斯引理）
\begin{sub}
\item 设 $f(x),g(x)\in\mathbb{Z}[x]$ 都是本原多项式。证明 $f(x)g(x)$ 也是本原多项式；
\item 设 $f(x)\in\mathbb{Z}[x]$ 本原，$\deg f\ge 1$。证明：若 $f(x)$ 在 $\mathbb{Q}[x]$ 上可约，
      则 $f(x)$ 在 $\mathbb{Z}[x]$ 上可分解为两个次数较低的本原整系数多项式之积。
\end{sub}

\item 设 $f_1(x),f_2(x),\ldots,f_s(x)\in P[x]$ 两两互素，$g_1(x),\ldots,g_s(x)\in P[x]$。证明：
\begin{sub}
\item 存在 $f(x)\in P[x]$，使得
\[
f(x)\equiv g_i(x)\pmod{f_i(x)},\qquad i=1,2,\ldots,s;
\]
\item 这样的 $f(x)$ 在模 $f_1f_2\cdots f_s$ 的意义下是唯一的。
\end{sub}

\item 设 $n\ge 2$，$A=(a_{ij})$ 为 $n$ 阶方阵，且 $A$ 的每一行元素之和与每一列元素之和都等于 $0$。
证明：$A$ 的所有代数余子式 $A_{ij}$ 都相等。
\end{qs}

\end{document}
